% 热木星（综述）
% license CCBYSA3
% type Wiki

本文根据 CC-BY-SA 协议转载翻译自维基百科\href{https://en.wikipedia.org/wiki/Hot_Jupiter}{相关文章}。

\begin{figure}[ht]
\centering
\includegraphics[width=6cm]{./figures/d9160f8109011126.png}
\caption{一幅艺术家绘制的近距离绕其母星运行的热木星示意图} \label{fig_HoJu_1}
\end{figure}
热木星（有时也称为热土星）是一类气态巨行星系外行星，被认为在物理性质上与木星相似（即木星类行星），但其轨道周期极短（(P<10) 天）$^{[1]}$。由于它们与母星距离极近、表面与大气温度很高，因而获得了“热木星”这一非正式名称$^{[2]}$。

在所有系外行星中，热木星是最容易通过径向速度法探测到的类型之一，这是因为它们在母星运动中引起的摆动幅度相对较大、变化也更快，明显强于其他已知类型的行星。其中最著名的热木星之一是 51 Pegasi b。它于 1995 年被发现，是第一颗被证实绕类太阳恒星运行的系外行星。51 Pegasi b 的轨道周期约为 4 天$^{[3]}$。
\subsection{一般特征}
\begin{figure}[ht]
\centering
\includegraphics[width=6cm]{./figures/e311dba254552faa.png}
\caption{截至 2014 年 1 月 2 日已发现的热木星（分布在左侧边缘，其中包括大多数通过凌日法探测到的行星，以黑点表示）} \label{fig_HoJu_2}
\end{figure}
\begin{figure}[ht]
\centering
\includegraphics[width=6cm]{./figures/c9637a190eec93c9.png}
\caption{具有隐藏水分子的热木星$^{[4]+}$} \label{fig_HoJu_3}
\end{figure}
尽管热木星之间存在多样性，但它们确实共享一些共同特征。
\begin{itemize}
\item 它们的定义性特征是较大的质量和极短的轨道周期，质量范围约为$0.36-11.8$个木星质量，轨道周期为$1.3-111$个地球日$^{[5]}$。其质量不能超过约 (13.6) 个木星质量，否则行星内部的压力和温度将高到足以引发氘核聚变，从而使该天体成为一颗褐矮星$^{[6]}$。
\item 大多数热木星具有近圆形轨道（低偏心率）。一般认为，其轨道会在邻近恒星的摄动或潮汐力作用下被逐渐圆化$^{[7]}$。它们是否能够在这种圆形轨道上长期稳定存在，或最终与其宿主恒星相撞，则取决于其轨道演化与物理演化之间的耦合关系，而这种耦合又通过能量耗散与潮汐形变联系在一起$^{[8]}$。
\item 许多热木星表现出异常低的平均密度。目前测得密度最低的是 TrES-4b，仅为$0.222,\mathrm{g/cm^3}$,$^{[9]}$。热木星的巨大半径尚未被完全理解，但通常认为其膨胀的大气包层可能源于强烈的恒星辐照、较高的大气不透明度、潜在的内部能量来源，以及与恒星距离极近导致行星外层超过洛希极限并被进一步拉伸$^{[9][10]}$。
\item 热木星通常处于潮汐锁定状态，即始终以同一面朝向其宿主恒星$^{[11]}$。
\item 由于其短轨道周期、相对较长的昼夜周期以及潮汐锁定特性，它们极可能拥有极端而奇特的大气环境$^{[3]}$。大气动力学模型预测，其大气将呈现显著的垂直分层结构，并伴随由辐射强迫以及热量与动量传输所驱动的强烈风场和超旋转的赤道急流$^{[12][13]}$。最新模型还预测存在多种风暴（涡旋），能够混合其大气并在不同区域之间输送冷热气体$^{[14]}$。
\item 在光球层尺度上，其昼夜温差被预测相当显著，例如基于 HD 209458 b 的模型给出的温差约为$500,\mathrm{K}$（约 $500,^\circ\mathrm{C}$，$900,^\circ\mathrm{F}$）$^{[13]}$。
\item 从统计上看，热木星更常见于 F 型和 G 型恒星周围，在 K 型恒星周围则较少，而围绕红矮星的热木星极为罕见$^{[15]}$。在对其分布进行概括时必须考虑各种观测偏差，但总体而言，其出现频率会随着恒星绝对星等的增加而呈指数式下降$^{[16]}$。
\end{itemize}